\documentclass[border=6pt,varwidth=14.5cm]{standalone}

\usepackage{algorithm2e}
\usepackage{amsmath}

\RestyleAlgo{plain}          % no float, no ruled box
\DontPrintSemicolon
\SetKwComment{Comment}{// }{}
\SetKwInOut{Input}{input}
\SetKwInOut{Output}{output}

% Notation follows Chapra, Applied Numerical Methods with MATLAB, Sec. 15.3
% (general linear least squares): design matrix Z and coefficient vector a.
% Course notation where it differs: the basis functions are phi_1..phi_m
% (z_j = phi_j(x)), vectors carry arrows, and theta is the generic name for
% the model parameters, here theta = a.

\begin{document}
\begin{minipage}{\linewidth}
\begin{algorithm}[H]
  \textbf{General Linear Least Squares}
  \BlankLine
  \Input{data $(\vec{x}_i, y_i)$ for $i = 1, \dots, n$; \
         basis functions $\phi_1, \dots, \phi_m$}
  \Output{$\theta^* = \vec{a} = (a_1, \dots, a_m)$, coefficients of
          $\hat{y} = a_1 \phi_1(\vec{x}) + \cdots + a_m \phi_m(\vec{x})$}
  \BlankLine
  \For(\tcp*[f]{build the $n \times m$ matrix $Z$}){$i \leftarrow 1$ \KwTo $n$}{
    \For{$j \leftarrow 1$ \KwTo $m$}{
      $Z_{ij} \leftarrow \phi_j(\vec{x}_i)$\;
    }
  }
  \vspace{3pt}
  solve $\big(Z^T Z\big)\, \vec{a} = Z^T \vec{y}$ for $\vec{a}$
    \tcp*{normal equations, $m \times m$}
  \vspace{3pt}
  \Return $\theta^* = \vec{a}$\;
\end{algorithm}
\end{minipage}
\end{document}
